\subsection{Regional persistence and uniqueness of the action selector}
\label{sec:accion-selector-regional}

The closure microfibre in Section~\ref{sec:generacion} supplies the
regional antecedent of the action reader. For the oriented ternary reduction
\(\Psi:\mathcal U_6\to\mathbb F_3^6\) already constructed, let
\(\mathcal F_\pi=\Psi^{-1}(010211)\). Its five emissions, separated
into initial and terminal blocks, are
\[
\begin{aligned}
 &(501,614,498\mid169,272,272),\\
 &(810,923,870\mid169,272,272),\\
 &(810,983,810\mid169,272,272),\\
 &(870,923,810\mid169,272,272),\\
 &(870,923,810\mid418,674,521).
\end{aligned}
\]
The return projection is
\[
 p_{\rm ret}:\mathbb Z^6\longrightarrow\mathbb Z^3,
 \qquad p_{\rm ret}(u_1,\ldots,u_6)=(u_4,u_5,u_6).
\]
Each emission is counted once at this step: persistence is selected
among the five regions, while the multiplicities of the seeds producing
those emissions remain a separate record.

\begin{proposition}[Persistent terminal block]
\label{prop:accion-bloque-persistente}
The multiset \(p_{\rm ret}(\mathcal F_\pi)\) has a unique mode,
\[
 b_\pi=(169,272,272),\qquad \operatorname{mult}(b_\pi)=4.
\]
The remaining block, \(b_\pi^-=(418,674,521)\), has multiplicity one.
\end{proposition}
\begin{proof}
The first four emissions project to \(b_\pi\); the fifth projects
to \(b_\pi^-\). The two blocks differ. The maximum multiplicity
is therefore four and is attained by exactly one block.
\end{proof}

The selector is defined on multisets of terminal blocks with a unique
mode of type \((r,s,s)\), \(r\ne s\). The group \(\mathfrak S_3\)
acts by relabelling the positions of these blocks; this equivariance
compares their representations without adding a dynamical symmetry to
the fixed fibre \(\Psi^{-1}(010211)\). The stabilizer
of \(b_\pi\) exchanges the positions occupied by \(272\) and fixes
the position occupied by \(169\). For a block of type \((r,s,s)\),
with \(r\ne s\), the value at this fixed position is
\[
 \operatorname{sing}(b)=\sum_{j=1}^3 b_j-2\operatorname{moda}(b).
\]
Selecting the regional mode and then applying this operation defines
the residual selector \(\rho_\pi\).

\begin{theorem}[Equivariant uniqueness of the return]
\label{thm:accion-selector-169}
Among selectors invariant under permutations of the five regions,
equivariant under permutations of terminal positions, selecting the
block of maximum multiplicity and retaining a position fixed by its
stabilizer, the return value is unique:
\[
 \rho_\pi=\operatorname{sing}(b_\pi)=169.
\]
\end{theorem}
\begin{proof}
Invariance with respect to the regions makes the selection depend on
the multiset rather than its enumeration. Modal persistence selects
\(b_\pi\) by Proposition~\ref{prop:accion-bloque-persistente}.
Its stabilizer has two position orbits: a pair and a fixed point.
The fixed-position requirement selects the latter, whose value is
\(169\). Under a permutation of positions, equivariance transports
this position while retaining its value. The construction satisfies
all four conditions, which also proves existence.
\end{proof}

The oriented difference between the two returns retains further data:
\[
 \omega_\pi=b_\pi^- - b_\pi=(249,402,249),\qquad
 \langle(1,1,1),\omega_\pi\rangle=900.
\]
The selector \(169\) is terminal data of degree zero. The difference
\(\omega_\pi\) assigns degree-one data to the oriented edge between
the two blocks; treating it as a cocycle requires specifying the complex
and checking the corresponding coboundary. This distinction preserves
both the return and its variation as their respective objects.

The selection is regional: in the full catalogue of \(468\) emissions,
\(169\) occurs \(84\) times, \(14\) in each position. Its role in
the reader follows jointly from the microfibre, modal persistence and
the stabilizer; its global frequency remains another count of the same
catalogue.

\subsection{Word length and the determinant-based return character}
\label{sec:accion-grado-determinante}

Let \(\mathscr P_{\rm pal}=\bigoplus_{n\ge0}\mathscr P_{{\rm pal},n}\) be the free
module on finite words over the emission alphabet, graded by length,
with the induced filtration. This domain contains regional emissions
and their concatenations; concatenation adds degrees. The formal length character
\[
 \chi_z([w])=z^{|w|},\qquad
 \chi_z([uv])=\chi_z([u])\chi_z([v])
\]
preserves this operation. Cylindrical compatibility of the reader sends
\(\mathscr P_{{\rm pal},n}\) to the homogeneous component \(z^n\mathbb R\) of
\(\mathbb R[z]\); evaluation at \(z=\alpha\) follows afterwards.
Since every \(w\in\mathcal F_\pi\) has length six, the residual
impulse is independent of the chosen region and has degree six:
\[
 \rho_\pi\chi_z([w])=169z^6,
 \qquad\left.\rho_\pi\chi_z([w])\right|_{z=\alpha}=169\alpha^6.
\]
The selection records one common impulse; it does not sum the five
emissions again after extracting their persistent return.
Emission length and path length in the incidence graph grade different
objects. In particular, closed paths through the four regional anchors
have minimum length eight; the degree used here counts the six positions
of an emission rather than the edges of such a path.

Dodecaphase closure and completion of the three nonadic pairs supply
the common coordinate \(A=10^3\alpha\). In radians,
\[
 x=\frac{\pi A}{180}=\frac{50\pi}{9}\alpha.
\]
Let \(R_{\rm h}\) be a real two-dimensional involution with zero trace,
\(R_{\rm h}^2=I_2\), \(\operatorname{tr}R_{\rm h}=0\). The two-sheet transport is
\[
 \mathcal T(x,y)=\exp(-xI_2-yR_{\rm h}).
\]
The oriented coordinate \(y\) remains unevaluated. Since the polynomial
\(X^2-1\) has distinct roots and \(R_{\rm h}\) has zero trace, its eigenvalues
are \(1,-1\). In an eigenbasis, the transport factors are
\(e^{-x-y}\) and \(e^{-x+y}\). Hence
\[
 \det\mathcal T=e^{-x-y}e^{-x+y}=e^{-2x}
 =\exp\!\left(-\frac{100\pi\alpha}{9}\right)=D_A.
\]
On the determinant line, \(\bigwedge^2\mathcal T\) acts as
\(D_A\) times the identity. Their product eliminates the oriented coordinate and preserves the
common contraction. The determinant is invariant under conjugation and
sheet exchange \(R_{\rm h}\mapsto-R_{\rm h}\); thus \(D_A\) is constructed before
the counter-angle is evaluated.

\begin{proposition}[Multiplicative return character]
\label{prop:accion-caracter-determinantal}
The multiplicative character \(\delta_A:(\mathbb N_0,+)\to(\mathbb R_{>0},\cdot)\)
whose value on one elementary extension is \(D_A\) is determined
by \(\delta_A(k)=D_A^k\).
\end{proposition}
\begin{proof}
A character preserves the identity, so \(\delta_A(0)=1\).
Concatenating \(k\) extensions gives
\(\delta_A(k)=\delta_A(1)^k=D_A^k\). These powers satisfy
multiplicativity and the prescribed elementary value, proving
existence and uniqueness.
\end{proof}

\Needspace{17\baselineskip}
\begin{definition}[Rules of regional continuation]
\label{def:accion-lector-regional}
The regional determinant-based reader preserves five operations:
\begin{enumerate}
 \item cylindrical length determines analytic degree;
 \item the return impulse and strict continuation are distinct additive
 terms in a renewal;
 \item after recording \(\rho_\pi\) once, continuation starts from
 the unit of the determinant line;
 \item each extension acts through \(\bigwedge^2\mathcal T\), with
 scalar factor \(D_A\);
 \item in the cardinal orientation of APP, regional return occupies the
 compensating sector and is subtracted from the fifth-degree action phase.
\end{enumerate}
\end{definition}

Unit normalization determines the amplitude of continuation as well as
its ratio. Indeed, all series
\[
 F_\lambda(z)=169z^6+\lambda\frac{D_Az^7}{1-D_Az}
 \qquad(\lambda\in\mathbb R)
\]
preserve the finite impulse and the same multiplicative recurrence for
tail coefficients. The third rule fixes \(\lambda=1\) before
evaluation at \(z=\alpha\). The residue is recorded once; extension
counts the continuing line with its unit. The renewal equation
\eqref{eq:revision-renovacion} and its uniqueness proof, developed below,
complete the chain from region and degree through transport to the
action section.
